\FigureLayoutDeclare{img/2000/syllabus_liberal/q21_fig1.png}{16.5cm}{46bp 0bp 231bp 0bp}
\chapter{2000年大纲卷（文）}
\section{单选题}

\begin{problem}
设集合 \(A = \{x \mid x \in \Z\text{ 且 }-10 \leqslant x \leqslant -1\}\)，\(B = \{x \mid x \in \Z\text{ 且 }|x| \leqslant 5\}\)，则 \(A \cup B\) 中的元素个数是（\quad）

\choices
    {\(11\)}
    {\(10\)}
    {\(16\)}
    {\(15\)}

\answer{C}

\begin{solution}
集合 \(A\) 含有整数 \(-10\)，\(-9\)，\(\ldots\)，\(-1\)，共有 \(10\) 个元素；集合 \(B\) 含有整数 \(-5\)，\(-4\)，\(\ldots\)，\(5\)，共有 \(11\) 个元素. 两集合的交集为 \(\{-5,-4,\ldots,-1\}\)，共有 \(5\) 个元素，所以
\[
|A\cup B|=10+11-5=16
\]
故选 C.
\end{solution}
\end{problem}

\begin{problem}
在复平面内，把复数 \(3 - \sqrt{3}\i\) 对应的向量按顺时针方向旋转 \(\frac{\pi}{3}\)，所得向量对应的复数是（\quad）

\choices
    {\(2\sqrt{3}\)}
    {\(-2\sqrt{3}\i\)}
    {\(\sqrt{3}-3\i\)}
    {\(3+\sqrt{3}\i\)}

\answer{B}

\begin{solution}
顺时针旋转 \(\frac{\pi}{3}\) 相当于乘以
\(\displaystyle \cos\frac{\pi}{3}-\i\sin\frac{\pi}{3}=\frac12-\frac{\sqrt3}{2}\i\). 因此所得复数为
\[
(3-\sqrt3\i)\left(\frac12-\frac{\sqrt3}{2}\i\right)=-2\sqrt3\i
\]
故选 B.
\end{solution}
\end{problem}

\begin{problem}
一个长方体共一顶点的三个面的面积分别是 \(\sqrt{2}\)，\(\sqrt{3}\)，\(\sqrt{6}\)，这个长方体对角线的长是（\quad）

\choices
    {\(2\sqrt{3}\)}
    {\(3\sqrt{2}\)}
    {\(6\)}
    {\(\sqrt{6}\)}

\answer{D}

\begin{solution}
设长方体的三条棱长分别为 \(x\)，\(y\)，\(z\)，则 \(xy=\sqrt2\)，\(yz=\sqrt3\)，\(zx=\sqrt6\). 由三式相乘得 \(xyz=\sqrt6\)，从而
\(x=\frac{xyz}{yz}=\sqrt2\)，\(y=\frac{xyz}{zx}=1\)，\(z=\frac{xyz}{xy}=\sqrt3\).
所以长方体的对角线长为 \(\sqrt{x^2+y^2+z^2}=\sqrt{2+1+3}=\sqrt6\). 故选 D.
\end{solution}
\end{problem}

\begin{problem}
已知 \(\sin\alpha>\sin\beta\)，那么下列命题成立的是（\quad）

\choices
    {若 \(\alpha\)，\(\beta\) 是第一象限角，则 \(\cos\alpha>\cos\beta\)}
    {若 \(\alpha\)，\(\beta\) 是第二象限角，则 \(\tan\alpha>\tan\beta\)}
    {若 \(\alpha\)，\(\beta\) 是第三象限角，则 \(\cos\alpha>\cos\beta\)}
    {若 \(\alpha\)，\(\beta\) 是第四象限角，则 \(\tan\alpha>\tan\beta\)}

\answer{D}

\begin{solution}
在第一象限内，正弦函数递增而余弦函数递减，所以 \(\sin\alpha>\sin\beta\) 推出 \(\cos\alpha<\cos\beta\)，命题 A 错. 第二象限内正切函数递增，而正弦函数递减，因此 \(\sin\alpha>\sin\beta\) 推出 \(\alpha<\beta\)，进而 \(\tan\alpha<\tan\beta\)，命题 B 错. 第三象限内正弦函数递减、余弦函数递增，故 \(\sin\alpha>\sin\beta\) 仍推出 \(\cos\alpha<\cos\beta\)，命题 C 错. 第四象限内正弦函数递增、正切函数递增，所以 \(\sin\alpha>\sin\beta\) 推出 \(\alpha>\beta\)，从而 \(\tan\alpha>\tan\beta\). 故选 D.
\end{solution}
\end{problem}

\begin{problem}
函数 \(y=-x\cos x\) 的部分图象是（\quad）

\choices
    {\choicebitmap[width=0.15\paperwidth]{img/2000/syllabus_liberal/q5_optA.png}}
    {\choicebitmap[width=0.15\paperwidth]{img/2000/syllabus_liberal/q5_optB.png}}
    {\choicebitmap[width=0.15\paperwidth]{img/2000/syllabus_liberal/q5_optC.png}}
    {\choicebitmap[width=0.15\paperwidth]{img/2000/syllabus_liberal/q5_optD.png}}

\answer{D}

\begin{solution}
函数 \(y=-x\cos x\) 是奇函数，因为 \(y(-x)=-y(x)\)，所以图象关于原点中心对称. 又 \(y(0)=0\)，当 \(x>0\) 足够小时 \(\cos x>0\)，于是 \(y<0\)；当 \(x<0\) 足够接近零时 \(y>0\). 此外，函数在 \(x=0\) 和 \(x=\frac{\pi}{2}\) 等处过 \(x\) 轴. 符合这些性质的图象是选项 D. 故选 D.
\end{solution}
\end{problem}

\begin{problem}
《中华人民共和国个人所得税法》规定，公民全月工资、薪金所得不超过 \(800\text{ 元}\) 的部分不必纳税，超过 \(800\text{ 元}\) 的部分为全月应纳税所得额，此项税款按下表分段累进计算：

\begin{center}
\begin{tabular}{|c|c|}
\hline
全月应纳税所得额 & 税率 \\
\hline
不超过 \(500\text{ 元}\) 的部分 & \(5\%\) \\
\hline
超过 \(500\text{ 元}\) 至 \(2000\text{ 元}\) 的部分 & \(10\%\) \\
\hline
超过 \(2000\text{ 元}\) 至 \(5000\text{ 元}\) 的部分 & \(15\%\) \\
\hline
\(\cdots\) & \(\cdots\) \\
\hline
\end{tabular}
\end{center}

某人一月份交纳此项税款 \(26.78\text{ 元}\)，则他的当月工资、薪金所得介于（\quad）

\choices
    {\(800\text{ 元}-900\text{ 元}\)}
    {\(900\text{ 元}-1200\text{ 元}\)}
    {\(1200\text{ 元}-1500\text{ 元}\)}
    {\(1500\text{ 元}-2800\text{ 元}\)}

\answer{C}

\begin{solution}
设应纳税所得额为 \(x\) 元. 前 \(500\) 元的税款为 \(500\times5\%=25\) 元，而总税款为 \(26.78\) 元，所以 \(x>500\). 在 \(500<x\leqslant2000\) 时，税款为
\[
25+10\%(x-500)=26.78
\]
解得 \(x=517.8\). 因此工资、薪金所得为 \(800+517.8=1317.8\) 元，介于 \(1200\) 元和 \(1500\) 元之间. 故选 C.
\end{solution}
\end{problem}

\begin{problem}
若 \(a>b>1\)，\(P=\sqrt{\lg a\cdot\lg b}\)，\(Q=\frac{1}{2}\left(\lg a+\lg b\right)\)，\(R=\lg\left(\frac{a+b}{2}\right)\)，则（\quad）

\choices
    {\(R<P<Q\)}
    {\(P<Q<R\)}
    {\(Q<P<R\)}
    {\(P<R<Q\)}

\answer{B}

\begin{solution}
令 \(u=\lg a\)，\(v=\lg b\)，则 \(u>v>0\). 由算术平均数与几何平均数的关系，\(P=\sqrt{uv}<\frac{u+v}{2}=Q\). 又因为 \(a>b>0\)，有 \(\frac{a+b}{2}>\sqrt{ab}\)，取常用对数得
\[
R=\lg\frac{a+b}{2}>\lg\sqrt{ab}=\frac12(\lg a+\lg b)=Q
\]
所以 \(P<Q<R\)，故选 B.
\end{solution}
\end{problem}

\begin{problem}
已知两条直线 \(l_{1}:y=x\)，\(l_{2}:ax-y=0\)，其中 \(a\) 为实数，当这两条直线的夹角在 \(\left(0,\frac{\pi}{12}\right)\) 内变动时，\(a\) 的取值范围是（\quad）

\choices
    {\((0,1)\)}
    {\(\left(\frac{\sqrt{3}}{3},\sqrt{3}\right)\)}
    {\(\left(\frac{\sqrt{3}}{3},1\right)\cup(1,\sqrt{3})\)}
    {\((1,\sqrt{3})\)}

\answer{C}

\begin{solution}
直线 \(l_1\) 的方向角为 \(\frac{\pi}{4}\). 若 \(a\leqslant0\)，则直线 \(l_2\) 与 \(l_1\) 的锐角夹角不小于 \(\frac{\pi}{4}\)，不可能属于 \(\left(0,\frac{\pi}{12}\right)\)，所以 \(a>0\). 令直线 \(l_2\) 的方向角为 \(\theta=\arctan a\)，则 \(0<\theta<\frac{\pi}{2}\). 由两直线夹角属于 \(\left(0,\frac{\pi}{12}\right)\)，有
\[
\frac{\pi}{6}<\theta<\frac{\pi}{3}
\]
于是
\[
\tan\frac{\pi}{6}<a<\tan\frac{\pi}{3}
\]
即 \(\frac{\sqrt3}{3}<a<\sqrt3\). 还要排除 \(a=1\)，因为此时两直线重合，夹角为 \(0\). 故
\[
a\in\left(\frac{\sqrt3}{3},1\right)\cup(1,\sqrt3)
\]
故选 C.
\end{solution}
\end{problem}

\begin{problem}
一个圆柱的侧面展开图是一个正方形，这个圆柱的全面积与侧面积的比是（\quad）

\choices
    {\(\frac{1+2\pi}{2\pi}\)}
    {\(\frac{1+4\pi}{4\pi}\)}
    {\(\frac{1+2\pi}{\pi}\)}
    {\(\frac{1+4\pi}{2\pi}\)}

\answer{A}

\begin{solution}
设圆柱的底面半径为 \(r\)，高为 \(h\). 侧面展开图为正方形，所以 \(h=2\pi r\). 圆柱的侧面积为 \(2\pi rh\)，全面积为 \(2\pi rh+2\pi r^2\)，故全面积与侧面积之比为
\[
\frac{2\pi rh+2\pi r^2}{2\pi rh}=1+\frac{r}{h}=1+\frac{1}{2\pi}=\frac{1+2\pi}{2\pi}
\]
故选 A.
\end{solution}
\end{problem}

\begin{problem}
过原点的直线与圆 \(x^{2}+y^{2}+4x+3=0\) 相切，若切点在第三象限，则该直线的方程是（\quad）

\choices
    {\(y=\sqrt{3}x\)}
    {\(y=-\sqrt{3}x\)}
    {\(y=\frac{\sqrt{3}}{3}x\)}
    {\(y=-\frac{\sqrt{3}}{3}x\)}

\answer{C}

\begin{solution}
圆的方程化为 \((x+2)^2+y^2=1\)，圆心为 \((-2,0)\)，半径为 \(1\). 过原点的切线可设为 \(y=mx\)，即 \(mx-y=0\). 由圆心到切线的距离等于半径，得
\[
\frac{|{-2m}|}{\sqrt{m^2+1}}=1
\]
所以 \(4m^2=m^2+1\)，即 \(m=\pm\frac{\sqrt3}{3}\). 切点在第三象限时，直线斜率必须为正，故 \(m=\frac{\sqrt3}{3}\). 因此切线方程为 \(y=\frac{\sqrt3}{3}x\)，故选 C.
\end{solution}
\end{problem}

\begin{problem}
过抛物线 \(y=ax^{2}\)（\(a>0\)）的焦点 \(F\) 作一直线交抛物线于 \(P\)，\(Q\) 两点，若线段 \(PF\) 与 \(FQ\) 的长分别是 \(p\)，\(q\)，则 \(\frac{1}{p}+\frac{1}{q}\) 等于（\quad）

\choices
    {\(2a\)}
    {\(\frac{1}{2a}\)}
    {\(4a\)}
    {\(\frac{4}{a}\)}

\answer{C}

\begin{solution}
将抛物线写成 \(x^2=4\rho y\) 的形式，其中 \(\rho=\frac{1}{4a}\)，可得焦点 \(F=\left(0,\frac{1}{4a}\right)\). 设过 \(F\) 的直线单位方向向量为 \((\cos\theta,\sin\theta)\)，用有向参数 \(s\) 表示直线上的点，则
\(x=s\cos\theta\)，\(y=\frac{1}{4a}+s\sin\theta\).
代入 \(y=ax^2\)，得
\[
4a^2\cos^2\theta\,s^2-4a\sin\theta\,s-1=0
\]
设两个根为 \(s_1\)，\(s_2\). 因为 \(P\)，\(Q\) 位于焦点两侧，所以 \(s_1s_2<0\)，并且
\(|s_1s_2|=\frac{1}{4a^2\cos^2\theta}\)，\(|s_1|+|s_2|=|s_1-s_2|=\frac{1}{a\cos^2\theta}\).
于是
\[
\frac1p+\frac1q=\frac{p+q}{pq}=\frac{1/(a\cos^2\theta)}{1/(4a^2\cos^2\theta)}=4a
\]
故选 C.
\end{solution}
\end{problem}

\begin{problem}
如图，\(OA\) 是圆锥底面中心 \(O\) 到母线的垂线，\(OA\) 绕轴旋转一周所得曲面将圆锥分成体积相等的两部分，则母线与轴的夹角为（\quad）

\bitmapfigure[max width=0.55\linewidth]{img/2000/syllabus_liberal/q12_fig1.png}

\choices
    {\(\arccos\frac{1}{\sqrt[3]{2}}\)}
    {\(\arccos\frac{1}{2}\)}
    {\(\arccos\frac{1}{\sqrt{2}}\)}
    {\(\arccos\frac{1}{\sqrt[4]{2}}\)}

\answer{D}

\begin{solution}
在过圆锥轴的截面内，设圆锥的高为 \(H\)，底面半径为 \(R\)，母线与轴的夹角为 \(\theta\). 以底面中心 \(O\) 为原点、轴为竖直方向，设顶点为 \(V=(0,H)\)，母线与底面圆周的交点为 \(B=(R,0)\)，垂足为 \(A=(r,h)\). 由于 \(A\) 在母线 \(VB\) 上，可设
\[
A=V+t(B-V)=\bigl(tR,H(1-t)\bigr)
\]
又因为 \(OA\perp VB\)，所以
\[
(tR,H(1-t))\cdot(R,-H)=0
\]
从而 \(t=\dfrac{H^2}{R^2+H^2}=\cos^2\theta\). 因此
\(\frac{r}{R}=\cos^2\theta\)，\(\frac{h}{H}=\sin^2\theta\).
旋转 \(OA\) 所得曲面将截面三角形分成两部分. 取靠近底面的部分，其体积等于半径为 \(R\)，\(r\)、高为 \(h\) 的圆台体积减去半径为 \(r\)、高为 \(h\) 的圆锥体积，即
\[
V_1=\frac{\pi h}{3}(R^2+Rr+r^2)-\frac{\pi r^2h}{3}=\frac{\pi h}{3}(R^2+Rr)
\]
题设说明 \(V_1\) 是原圆锥体积的一半，因此
\[
\frac{\pi h}{3}(R^2+Rr)=\frac12\cdot\frac{\pi R^2H}{3}
\]
令 \(u=\cos^2\theta=r/R\)，则 \(h/H=1-u\)，从而 \(2(1-u)(1+u)=1\)，即 \(u^2=\frac12\). 因为 \(0<\theta<\frac{\pi}{2}\)，所以 \(u=\frac1{\sqrt2}\)，故
\[
\cos\theta=\frac1{\sqrt[4]{2}}
\]
故选 D.
\end{solution}
\end{problem}

\section{填空题}

\begin{problem}
乒乓球队的 \(10\) 名队员中有 \(3\) 名主力队员，派 \(5\) 名参加比赛，\(3\) 名主力队员要安排在第一、三、五位置，其余 \(7\) 名队员选 \(2\) 名安排在第二、四位置，那么不同的出场安排共有\fillinblank{}种（用数字作答）.

\answer{252}

\begin{solution}
三名主力队员分别安排在第一、三、五位置，有 \(3!=6\) 种安排. 其余两名队员从七名非主力队员中选出并安排在第二、四位置，有 \(\mathrm A_7^2=7\times6=42\) 种安排. 因此总数为
\[
3!\mathrm A_7^2=6\times42=252
\]
\end{solution}
\end{problem}

\begin{problem}
椭圆 \(\frac{x^{2}}{9}+\frac{y^{2}}{4}=1\) 的焦点 \(F_{1}\)，\(F_{2}\)，点 \(P\) 为其上的动点，当 \(\angle F_{1}PF_{2}\) 为钝角时，点 \(P\) 横坐标的取值范围是\fillinblank{}.

\answer{\(\left(-\frac{3}{\sqrt{5}},\frac{3}{\sqrt{5}}\right)\)}

\begin{solution}
椭圆的半长轴、半短轴分别为 \(3,2\)，所以焦距参数满足 \(c=\sqrt{3^2-2^2}=\sqrt5\)，可取 \(F_1=(-\sqrt5,0)\)，\(F_2=(\sqrt5,0)\). 设 \(P=(x,y)\)，则
\[
\overrightarrow{PF_1}\cdot\overrightarrow{PF_2}=x^2+y^2-5
\]
角 \(F_1PF_2\) 为钝角当且仅当该数量积小于 \(0\)，即 \(x^2+y^2<5\). 又 \(P\) 在椭圆上，故 \(y^2=4(1-x^2/9)\)，从而
\[
x^2+y^2=4+\frac59x^2<5
\]
因此 \(x^2<\frac95\)，即
\[
x\in\left(-\frac3{\sqrt5},\frac3{\sqrt5}\right)
\]
\end{solution}
\end{problem}

\begin{problem}
设 \(\{a_{n}\}\) 是首项为 \(1\) 的正项数列，且 \((n+1)a_{n+1}^{2}-na_{n}^{2}+a_{n+1}a_{n}=0\)（\(n=1\)，\(2\)，\(3\)，\(\dots\)），则它的通项公式是 \(a_{n}=\fillinblank{}\).

\answer{\(\frac{1}{n}\)}

\begin{solution}
由题设，递推式左边可以因式分解为
\[
\bigl((n+1)a_{n+1}-na_n\bigr)(a_{n+1}+a_n)=0
\]
数列为正项数列，所以 \(a_{n+1}+a_n>0\)，从而 \((n+1)a_{n+1}=na_n\). 于是 \(na_n\) 为常数，结合 \(a_1=1\) 得
\(na_n=a_1=1\)，\(a_n=\frac1n\).
\end{solution}
\end{problem}

\begin{problem}
如图，\(E\)，\(F\) 分别为正方体面 \(ADD_{1}A_{1}\)，面 \(BCC_{1}B_{1}\) 的中心，则四边形 \(BFD_{1}E\) 在该正方体的面上的射影可能是\fillinblank{}（要求：把可能的图序号都填上）.

\examfiguregroup[float=inline]{img/2000/syllabus_liberal/q16_fig1.png}{%
    \vbox{%
        \hbox to \linewidth{%
            \hfil\bitmapinclude[max width=0.82\linewidth]{img/2000/syllabus_liberal/q16_fig1.png}\hfil
        }%
        \vskip 0.4\baselineskip
        \hbox to \linewidth{%
            \hfil\bitmapinclude[width=0.20\linewidth]{img/2000/syllabus_liberal/q16_optA.png}\hfil
            \bitmapinclude[width=0.20\linewidth]{img/2000/syllabus_liberal/q16_optB.png}\hfil
            \bitmapinclude[width=0.20\linewidth]{img/2000/syllabus_liberal/q16_optC.png}\hfil
            \bitmapinclude[width=0.20\linewidth]{img/2000/syllabus_liberal/q16_optD.png}\hfil
        }%
    }%
}

\answer{\(\circled{2}\)，\(\circled{3}\)}

\begin{solution}
取正方体棱长为 \(1\)，建立坐标系 \(A=(0,0,0)\)，\(B=(1,0,0)\)，\(D=(0,1,0)\)，\(D_1=(0,1,1)\). 由 \(E\)，\(F\) 分别为两个相对侧面的中心，得
\(E=\left(0,\frac12,\frac12\right)\)，\(F=\left(1,\frac12,\frac12\right)\).
四边形 \(BFD_1E\) 的两组对边分别平行，所以它本身是平行四边形. 正投影保持平行关系. 投影到 \(x=0\) 或 \(x=1\) 的面时，四个顶点的投影共线，四边形退化为线段，对应图 \(\circled{3}\). 投影到其余四个面，即 \(y=0\)、\(y=1\)、\(z=0\)、\(z=1\) 所在的面时，四个顶点的投影均不重合，仍得到非退化平行四边形，对应图 \(\circled{2}\). 因此可能的图形是平行四边形和线段，即图序号 \(\circled{2}\)、\(\circled{3}\).
\end{solution}
\end{problem}

\section{解答题}

\begin{problem}
已知函数 \(y=\sqrt{3}\sin x+\cos x\)（\(x\in\R\)）.

\begin{enumerate}
\item 当函数 \(y\) 取得最大值时，求自变量 \(x\) 的集合.
\item 该函数的图象可由 \(y=\sin x\)（\(x\in\R\)）的图象经过怎样的平移和伸缩变换得到？
\end{enumerate}

\begin{solution}
\begin{enumerate}
\item \(y=\sqrt3\sin x+\cos x=2\sin\left(x+\frac{\pi}{6}\right)\). 函数取得最大值 \(2\) 当且仅当 \(x+\frac{\pi}{6}=\frac{\pi}{2}+2k\pi\)，所以 \(x=\frac{\pi}{3}+2k\pi\)，其中 \(k\in\Z\).
\item 由 \(y=2\sin\left(x+\frac{\pi}{6}\right)\) 可知，先将 \(y=\sin x\) 的图象向左平移 \(\frac{\pi}{6}\)，再沿纵轴方向伸长为原来的 \(2\) 倍即可.
\end{enumerate}
\end{solution}
\end{problem}

\begin{problem}
设 \(\{a_{n}\}\) 为等差数列，\(S_{n}\) 为数列 \(\{a_{n}\}\) 的前 \(n\) 项和，已知 \(S_{7}=7\)，\(S_{15}=75\)，\(T_{n}\) 为数列 \(\left\{\frac{S_{n}}{n}\right\}\) 的前 \(n\) 项和，求 \(T_{n}\).

\begin{solution}
设等差数列首项为 \(a_1\)，公差为 \(d\). 由 \(S_n=\frac n2(2a_1+(n-1)d)\) 及已知条件，得
\(2a_1+6d=2\)，\(2a_1+14d=10\).
解得 \(d=1\)，\(a_1=-2\). 因此
\(S_n=\frac n2(n-5)\)，\(\frac{S_n}{n}=\frac{n-5}{2}\).
于是
\[
T_n=\sum_{k=1}^n\frac{k-5}{2}=\frac14n(n+1)-\frac52n=\frac14n^2-\frac94n
\]
\end{solution}
\end{problem}

\begin{problem}
如图，已知平行六面体 \(ABCD-A_{1}B_{1}C_{1}D_{1}\) 的底面 \(ABCD\) 为菱形，且 \(\angle C_{1}CB=\angle C_{1}CD=\angle BCD\).

\begin{enumerate}
\item 证明：\(C_{1}C\perp BD\)；
\item 当 \(\frac{CD}{CC_{1}}\) 的值为多少时，能使 \(A_{1}C\perp\) 平面 \(C_{1}BD\)？请给出证明.
\end{enumerate}

\bitmapfigure[float=inline,max width=0.7\linewidth]{img/2000/syllabus_liberal/q19_fig1.png}

\begin{solution}
\begin{enumerate}
\item 设 \(\bs u=\overrightarrow{CB}\)，\(\bs v=\overrightarrow{CD}\)，\(\bs w=\overrightarrow{CC_1}\)，并记 \(\lvert\bs u\rvert=\lvert\bs v\rvert=s\)，\(\angle BCD=\varphi\). 由底面为菱形，\(\lvert\bs u\rvert=\lvert\bs v\rvert=s\). 又由题设三个角相等，
\[
\bs w\cdot\bs u=\lvert\bs w\rvert s\cos\varphi=\bs w\cdot\bs v
\]
所以 \(\bs w\cdot(\bs v-\bs u)=0\)，而 \(\bs v-\bs u=\overrightarrow{BD}\)，故 \(C_1C\perp BD\).
\item 有 \(\overrightarrow{CA_1}=\bs u+\bs v+\bs w\). 由（1）立即得到
\[
(\bs u+\bs v+\bs w)\cdot(\bs v-\bs u)=0
\]
再取平面 \(C_1BD\) 内的方向向量 \(\overrightarrow{BC_1}=\bs w-\bs u\)，则
\[
(\bs u+\bs v+\bs w)\cdot(\bs w-\bs u)
=\lvert\bs w\rvert^2-s^2+s\lvert\bs w\rvert\cos\varphi-s^2\cos\varphi
=\bigl(\lvert\bs w\rvert-s\bigr)\bigl(\lvert\bs w\rvert+s+s\cos\varphi\bigr)
\]
因 \(0<\varphi<\pi\)，第二个因子为正，所以垂直条件等价于 \(\lvert\bs w\rvert=s\). 因此
\[
\frac{CD}{CC_1}=\frac{s}{\lvert\bs w\rvert}=1
\]
\end{enumerate}
\end{solution}
\end{problem}

\begin{problem}
设函数 \(f(x)=\sqrt{x^{2}+1}-ax\)，其中 \(a>0\).

\begin{enumerate}
\item 解不等式 \(f(x)\leqslant1\)；
\item 证明：当 \(a\geqslant1\) 时，函数 \(f(x)\) 在区间 \([0,+\infty)\) 上是单调函数.
\end{enumerate}

\begin{solution}
\begin{enumerate}
\item 当 \(x<0\) 时，\(\sqrt{x^2+1}>1\)，且 \(-ax>0\)，所以 \(f(x)>1\)，不等式无负数解. 当 \(x\geqslant0\) 时，右端 \(1+ax>0\)，故原不等式等价于
\[
x^2+1\leqslant(1+ax)^2
\]
即 \(x((1-a^2)x-2a)\leqslant0\). 若 \(0<a<1\)，得到 \(0\leqslant x\leqslant\frac{2a}{1-a^2}\)；若 \(a\geqslant1\)，对所有 \(x\geqslant0\) 不等式都成立. 因此解集为
\[
\begin{cases}
\left[0,\dfrac{2a}{1-a^2}\right],&0<a<1,\\
[0,+\infty),&a\geqslant1.
\end{cases}
\]
\item 取 \(0\leqslant x_1<x_2\). 有
\[
f(x_2)-f(x_1)=(x_2-x_1)\left(\frac{x_1+x_2}{\sqrt{x_2^2+1}+\sqrt{x_1^2+1}}-a\right)
\]
由于 \(x_i\geqslant0\)，分式小于 \(1\)，而 \(a\geqslant1\)，故括号内小于 \(0\). 于是 \(f(x_2)<f(x_1)\)，所以 \(f(x)\) 在 \([0,+\infty)\) 上单调递减.
\end{enumerate}
\end{solution}
\end{problem}

\begin{problem}
某蔬菜基地种植西红柿，由历年市场行情得知，从二月一日起的 \(300\) 天内，西红柿市场售价与上市时间的关系用图一的一条折线表示；西红柿的种植成本与上市时间的关系用图二的抛物线段表示.

\begin{enumerate}
\item 写出图一表示的市场售价与时间的函数关系式 \(p=f(t)\)；写出图二表示的种植成本与时间的函数关系式 \(Q=g(t)\)；
\item 认定市场售价减去种植成本为纯收益，问何时上市的西红柿纯收益最大？
\end{enumerate}

注：市场售价与种植成本的单位：\(\text{元}/10^{2}\mathrm{~kg}\)，时间单位：\(\text{天}\).

\begingroup
\leftskip=0pt\relax
\everypar{}
\bitmapfigure[float=inline,max width=0.9\linewidth]{img/2000/syllabus_liberal/q21_fig1.png}
\endgroup

\begin{solution}
\begin{enumerate}
\item 由图一中折线经过 \((0,300)\)、\((200,100)\)、\((300,300)\)，所以
\[
p=f(t)=\begin{cases}
300-t, & 0\leqslant t\leqslant200,\\
2t-300, & 200<t\leqslant300.
\end{cases}
\]
由图二的顶点 \((150,100)\) 及点 \((50,150)\) 可得
\(Q=g(t)=\frac1{200}(t-150)^2+100\)，\(0\leqslant t\leqslant300\).
\item 设纯收益为 \(w(t)=p(t)-g(t)\). 当 \(0\leqslant t\leqslant200\) 时，
\[
w(t)=-\frac1{200}t^2+\frac12t+\frac{175}{2}
\]
其对称轴为 \(t=50\)，故这一段的最大值在 \(t=50\) 处取得，且为 \(100\). 当 \(200<t\leqslant300\) 时，
\[
w(t)=-\frac1{200}t^2+\frac72t-\frac{1025}{2}
\]
其对称轴为 \(t=350\)，所以在区间 \((200,300]\) 上递增，最大值为 \(w(300)=\frac{175}{2}<100\). 因此从二月一日起 \(50\) 天时上市，纯收益最大.
\end{enumerate}
\end{solution}
\end{problem}

\begin{problem}
如图，已知梯形 \(ABCD\) 中 \(|AB|=2|CD|\)，点 \(E\) 分有向线段 \(\overrightarrow{AC}\) 所成的比为 \(\lambda\)，双曲线过 \(C\)，\(D\)，\(E\) 三点，且以 \(A\)，\(B\) 为焦点，当 \(\frac{2}{3}\leqslant\lambda\leqslant\frac{3}{4}\) 时，求双曲线离心率 \(e\) 的取值范围.

\bitmapfigure[max width=0.55\linewidth]{img/2000/syllabus_liberal/q22_fig1.png}

\begin{solution}
以 \(AB\) 的垂直平分线为 \(y\) 轴、\(AB\) 所在直线为 \(x\) 轴建立直角坐标系. 设双曲线的焦点为 \(A=(-c,0)\)、\(B=(c,0)\)，方程为
\(\frac{x^2}{a^2}-\frac{y^2}{b^2}=1\)，\(c^2=a^2+b^2\).
因为 \(CD\parallel AB\)，且 \(C\)，\(D\) 在同一双曲线上，关于 \(y\) 轴对称. 又 \(AB=2CD\)，故可设
\(C=\left(\frac c2,h\right)\)，\(D=\left(-\frac c2,h\right)\).
由点 \(C\) 在双曲线上，得
\(\frac{c^2}{4a^2}-\frac{h^2}{b^2}=1\)，\(\frac{h^2}{b^2}=\frac14\frac{c^2}{a^2}-1\).
按有向线段定比分点的定义，\(\overrightarrow{AE}=\lambda\overrightarrow{EC}\)，所以
\[
E=\frac{A+\lambda C}{1+\lambda}
=\left(\frac{c(\lambda-2)}{2(1+\lambda)},\frac{\lambda h}{1+\lambda}\right)
\]
令 \(e^2=\frac{c^2}{a^2}\)，把点 \(E\) 代入双曲线方程，并利用上式，得到
\[
\frac{e^2(\lambda-2)^2}{4(1+\lambda)^2}
-\frac{\lambda^2}{(1+\lambda)^2}\left(\frac{e^2}{4}-1\right)=1
\]
整理得 \(e^2(1-\lambda)=1+2\lambda\)，即
\[
e^2=\frac{1+2\lambda}{1-\lambda}
\]
函数 \(\frac{1+2\lambda}{1-\lambda}\) 在 \(\left[\frac23,\frac34\right]\) 上递增，因此
\[
7\leqslant e^2\leqslant10
\]
双曲线离心率为正，故
\[
e\in[\sqrt7,\sqrt{10}]
\]
\end{solution}
\end{problem}
